\documentclass[aspectratio=169]{beamer}

\usepackage[utf8]{inputenc}
\usepackage[T1]{fontenc}
\usepackage[polish]{babel}
\usepackage{lmodern}
\usepackage{graphicx}
\usepackage{amssymb}
% \usepackage{amsthm} % ZAKOMENTOWANE: Powoduje błąd w Beamerze (konflikt z definicją 'proof')
\usepackage{algorithm}
\usepackage{algorithmic}
\usepackage{booktabs}
\usepackage{caption} % Dodano dla captionof
\usepackage{hyperref}
% Pakiet cite opcjonalnie dla lepszego formatowania numerów
% \usepackage{cite}

% Konfiguracja linków
\hypersetup{
    colorlinks=true,
    linkcolor=blue,
    urlcolor=blue,
    citecolor=blue
}

% Motyw
\usetheme{Madrid}
\usecolortheme{beaver}

% Dane prezentacji
\title[RL w modelach z dyskontowaniem QH]{Algorytmy oparte o~wzmocnione nauczanie maszynowe w~dyskontowanych modelach markowskich}
\author{Michał Wrona}
\institute[PWr]{Politechnika Wrocławska\\Wydział Matematyki\\[0.5em]Promotor: prof. dr hab. inż. Anna Jaśkiewicz}
\date{\today}

% Spolszczenie etykiet algorytmów
\floatname{algorithm}{Algorytm}
\renewcommand{\listalgorithmname}{Spis algorytmów}
\renewcommand{\algorithmicrequire}{\textbf{Wymagania:}}
\renewcommand{\algorithmicensure}{\textbf{Wynik:}}

% Automatyczny spis treści przed każdą sekcją
\AtBeginSection[]{
    \begin{frame}
    \frametitle{Plan prezentacji}
    \tableofcontents[currentsection]
    \end{frame}
}

\begin{document}

% Strona tytułowa
\begin{frame}[plain]
    \titlepage
\end{frame}

\begin{frame}{Plan prezentacji}
    \tableofcontents
\end{frame}

% ==========================================
% Wstęp
% ==========================================
\section{Wprowadzenie}

\begin{frame}{Motywacja i~cel pracy}
\textbf{Motywacja:}    
\begin{itemize}
    \item \textbf{Uczenie maszynowe poprzez wzmocnienie (RL)}
    \vspace{0.3em}
    \item \textbf{Połączenie psychologii z~matematyką}
    \begin{itemize}
        \item Ekonomia behawioralna
        \item Zjawisko prokrastynacji
    \end{itemize}
\end{itemize}
\textbf{Cel Pracy:}    
\begin{itemize}
    \item Opracowanie i analiza algorytmów obliczania funkcji wartości
    oraz optymalnej polityki dla decydenta zobowiązanego w modelach z dyskontowaniem quasi-hiperbolicznym
     oraz udowodnienie ich poprawności na prostych przykładach możliwych do policzenia
     analitycznie.
       
\end{itemize}
\end{frame}

\begin{frame}{Spis symboli i oznaczeń}
    \scriptsize % Zmniejszenie czcionki, aby zmieścić listę na jednym slajdzie
    \begin{columns}[T] % T = wyrównanie kolumn do góry
        % Lewa kolumna
        \begin{column}{0.48\textwidth}
            \begin{itemize}
                \setlength{\itemsep}{1pt} % Zmniejszenie odstępów
                \item[$M$] -- markowski proces decyzyjny
                \item[$S$] -- zbiór stanów
                \item[$A$] -- zbiór akcji
                \item[$q$] -- prawdopodobieństwo przejścia
                \item[$r$] -- funkcja nagrody
                \item[$d(t)$] -- funkcja dyskontująca
                \item[$\alpha$] -- parametr uprzedzenia teraźniejszości
                \item[$\beta$] -- współczynnik dyskontowania
                \item[$\Pr(A)$] -- zbiór miar probabilistycznych na $A$
                \item[$\mathcal{H}_n$] -- historie długości $n$
                \item[$\Phi$] -- zbiór polityk zależnych od historii
                \item[$\Phi_m$] -- zbiór polityk markowskich
                \item[$\Phi_s$] -- zbiór polityk stacjonarnych
                \item[$\phi$] -- polityka
            \end{itemize}
        \end{column}

        % Prawa kolumna
        \begin{column}{0.48\textwidth}
            \begin{itemize}
                \setlength{\itemsep}{1pt}
                \item[$\phi_s$] -- polityka stacjonarna
                \item[$\mu$] -- reguła pierwszego kroku
                \item[$\mu^\star$] -- optymalna reguła pierwszego kroku
                \item[$\phi_s^\star$] -- optymalna polityka stacjonarna
                \item[$\phi_\beta^\star$] -- optymalna polityka wykładnicza
                \item[$J_{\phi}^{\alpha,\beta}$] -- wartość quasi-hiperboliczna
                \item[$J_{\phi}^{\beta}$] -- wartość wykładnicza
                \item[$J_n, W_n$] -- przybliżenia iteracyjne ($Q_n$ dla par)
                \item[$\rho_n^J, \rho_n^W$] -- wagi wiarygodności
                \item[$\eta_n, \theta_n$] -- kroki uczenia
                \item[$\nu$] -- polityka próbkowania
                \item[$T^{\phi_s}$] -- operator Bellmana
                \item[$\|\cdot\|_\infty$] -- norma maksimum
                \item[$\mathbb{E}$] -- wartość oczekiwana
            \end{itemize}
        \end{column}
    \end{columns}
\end{frame}

\begin{frame}{Markowski proces decyzyjny (MDP)}
\textbf{Model:} $M = (S, A, q, r, \alpha, \beta)$
\begin{columns}[T]
\begin{column}{0.48\textwidth}
\begin{itemize}
    \item $S$ -- skończony zbiór stanów
    \item $A$ -- skończony zbiór akcji
    \item $q(s' \mid s, a)$ -- prawdopodobieństwo przejścia
    \item $r(s, a)$ -- funkcja nagrody
    \item $\alpha > 0$ -- parametr uprzedzenia teraźniejszości
    \item $\beta \in [0,1)$ -- współczynnik dyskontowania
\end{itemize}
\end{column}
\begin{column}{0.48\textwidth}
\centering
\includegraphics[width=0.95\textwidth]{schematRL.png}
\small Schemat interakcji agent--środowisko
\end{column}
\end{columns}

\vspace{0.5em}
\textbf{Funkcja dyskontująca:}
\[
d(t) = 
\begin{cases}
    1, & \text{jeśli } t = 0, \\
    \alpha\,\beta^{t}, & \text{jeśli } t \ge 1.
\end{cases}
\]
\end{frame}

\begin{frame}{Polityki i~funkcje wartości}
\textbf{Typy polityk:}
\begin{itemize}
    \item $\Phi$ -- polityki zależne od historii: $\phi_n : \mathcal{H}_n \rightarrow \Pr(A)$
    \item $\Phi_m$ -- polityki markowskie: $\phi_n : S \rightarrow \Pr(A)$
    \item $\Phi_s$ -- polityki stacjonarne: $\phi_s : S \rightarrow \Pr(A)$, $\phi_n = \phi_s$ dla każdego $n$
\end{itemize}

\vspace{0.5em}
\textbf{Funkcja wartości z~dyskontowaniem QH:}
\begin{align*}
    J_{\phi}^{\alpha,\beta}(s_0)
    &= \mathbb{E}\!\left[ r(s_0,a_0) + \alpha\,\beta\, J_{\overrightarrow{\phi}}^{\beta}(s_1) \right] \\
    &= \sum_a \phi_0(a \mid s_0) \left[ r(s_0,a) + \alpha\,\beta\, \sum_{s'} q(s' \mid s_0,a) J_{\overrightarrow{\phi}}^{\beta}(s') \right]
\end{align*}
gdzie $\overrightarrow{\phi} = (\phi_1, \phi_2, \ldots)$, a~$J_{\overrightarrow{\phi}}^{\beta}$ to wartość z~dyskontowaniem wykładniczym.
\end{frame}

\begin{frame}{Typy decydentów w~zależności od $\alpha$}
\begin{columns}[T]
\begin{column}{0.5\textwidth}
\begin{itemize}
    \item $\alpha = 1$ -- \textbf{decydent standardowy} (czasowo spójny)
    \item $0 < \alpha < 1$ -- \textbf{decydent quasi-hiperboliczny} (czasowo niespójny)
    \begin{itemize}
        \item \textit{Naiwny} -- nie przewiduje własnych przyszłych preferencji
        \item \textit{Zobowiązany} (precommitted) -- trzyma się wcześniej ustalonego planu
        \item \textit{Wyrafinowany} (sophisticated) -- przewiduje i~bierze pod uwagę przyszłe zmiany preferencji
    \end{itemize}
    \item $\alpha = 0$ -- decydent krótkowzroczny
    \item $\alpha > 1$ -- decydent długowzroczny
\end{itemize}
\vspace{0.5em}
W pracy koncentrujemy się na decydentach zobowiązanych (precommitted agents).
\end{column}
\begin{column}{0.42\textwidth}
\centering
\includegraphics[width=0.95\textwidth]{qh_discounting_comparison.png}
\small Porównanie funkcji dyskontowania dla różnych wartości $\alpha$ przy $\beta = 0.95$
\end{column}
\end{columns}
\end{frame}

% --- SLAJD 1: MODEL ---
\begin{frame}{Model decyzji zakupowej}
    \begin{columns}[T]
        \begin{column}{0.6\textwidth}
            \textbf{Problem:} Impulsowy zakup ($t=0$) vs poszukiwanie okazji ($t>0$).
            
            \vspace{0.5em}
            \textbf{Parametry wejściowe (skala 1--5):}
            \begin{itemize}
                \item Cena ($p$), Majątek ($w$), Wiek ($a$).
            \end{itemize}
    
            \vspace{0.5em}
            \textbf{Wyznaczenie cierpliwości $\alpha$:}
            \begin{itemize}
                \item Normalizacja zmiennych do $[0,1]$.
                \item $\text{score} = \frac{1}{3}(p_{\text{norm}} + w_{\text{norm}} + a_{\text{norm}})$.
                \item $\alpha = \sigma(\text{score})$ (przesunięta sigmoida).
            \end{itemize}
            
            \vspace{0.5em}
            \textbf{Funkcja nagrody:}
            \[ R(t) = 1 + \text{bonus} \cdot (1 - e^{-0.5t}) \]
            Decydent maksymalizuje $u_t = \alpha \beta^t R(t)$.
        \end{column}
        
        \begin{column}{0.38\textwidth}
            \centering
            \includegraphics[width=\textwidth]{Sigmoid_alpha.png}
            \vspace{-0.5em}
            \captionof{figure}{\scriptsize Zależność $\alpha$ od wyniku \textit{score}.}
            \label{fig:sigmoid-alpha}
            \vspace{0.5em}
            {\scriptsize $\sigma(x; k, x_0) = \tfrac{1}{2(1 + e^{-k(x - x_0)})} + 0.5$}
        \end{column}
    \end{columns}
\end{frame}

% --- SLAJD 2: WYNIKI ---
\begin{frame}{Przykład 2: Wyniki symulacji dla profili klientów}
    \centering
    % Wykres procesu decyzyjnego
    \includegraphics[width=0.85\textwidth]{qhExp2.png}
    
    \vspace{0.5em}
    
    % Tabela z wynikami
    \resizebox{0.95\textwidth}{!}{%
    \begin{tabular}{l|ccc|cc}
    \toprule
    \textbf{Profil} & \textbf{Cena} & \textbf{Majątek} & \textbf{Wiek} & \textbf{$\alpha$} & \textbf{Decyzja ($t^\star$)} \\ \midrule
    Młody Bogacz (tani prod.) & 2 (tani) & 1 (b. bogaty) & 2 (młody) & 0.544 & Kup teraz (0) \\
    Emeryt (drogi prod.) & 4 (drogi) & 4 (biedny) & 5 (b. stary) & 0.956 & Przeszukaj (6 kr.) \\
    Student (prod. norm.) & 3 (norm.) & 4 (biedny) & 2 (młody) & 0.750 & Kup teraz (0) \\ \bottomrule
    \end{tabular}%
    }
\end{frame}

% ==========================================
% Teoria i algorytmy
% ==========================================
\section{Teoria i~algorytmy}

\begin{frame}{Twierdzenie o~redukcji polityk}
\begin{block}{Twierdzenie 1 \cite{eshwar2025}}
Dla MDP o~skończonych zbiorach stanów i~akcji, dla dowolnej polityki $\phi = (\phi_0, \phi_1, \ldots)$ istnieje reguła pierwszego kroku $\mu$ oraz stacjonarna polityka $\phi_s$ takie, że:
\[
J_{\phi}^{\alpha,\beta}(s) = J_{(\mu,\phi_s)}^{\alpha,\beta}(s), \qquad \forall s \in S.
\]
\end{block}

\vspace{1em}
\textbf{Dowód:} 1285\_MS.pdf.
\end{frame}


\begin{frame}{Operator Bellmana i~kontrakcja}
\begin{block}{Lemat o~kontrakcji \cite{puterman2014,Kallenberg2022}}
Operator $T^{\phi_s}: \mathbb{R}^{|S|} \to \mathbb{R}^{|S|}$ zdefiniowany jako:
\[
(T^{\phi_s} J)(s) = \sum_a \phi_s(a\mid s) \left[ r(s, a) + \sum_{s'} q(s'|s, a) \left[ -(1 - \alpha)\beta \mathbb{E}_{\phi_s}[r(s', \cdot)] + \beta J(s') \right] \right]
\]
jest kontrakcją w~normie $\|\cdot\|_\infty$ ze współczynnikiem $\beta < 1$.
\end{block}

\vspace{1em}
\textbf{Konsekwencje:} Istnieje jedyny punkt stały $J^{\alpha, \beta}_{\phi_s}$ oraz iteracja $J_{n+1} = T^{\phi_s}(J_n)$ zbiega wykładniczo.

\vspace{0.5em}
\textbf{Dowód:} 1285\_MS.pdf.
\end{frame}

\begin{frame}{Algorytm 1: Obliczenie funkcji wartości dla dyskontowania QH}
\begin{algorithm}[H]
\caption{Obliczenie funkcji wartości dla dyskontowania QH \cite{eshwar2025}}
\label{alg:policy-eval-defense}
\begin{algorithmic}[1]
\small
\STATE \textbf{Wymagania:} Kroki uczenia $(\eta_n)$, $(\theta_n)$; współczynniki dyskontowania $\alpha$, $\beta$; polityka próbkowania $\nu$; polityki oceniane $\mu$, $\phi_s$
\STATE \textbf{Inicjalizacja:} $J_0, W_0 \in \mathbb{R}^{|S|}$
\FOR{$n = 0, 1, 2, \dots$}
\FOR{$s \in S$}
\STATE Próbkuj $a \sim \nu(\cdot|s)$
\STATE Zaobserwuj nagrodę $r(s, a)$ oraz stan kolejny $s' \sim q(\cdot|s, a)$
\STATE Próbkuj $a' \sim \phi_s(\cdot|s')$
\STATE Zaobserwuj nagrodę $r(s', a')$
\STATE Oblicz cel TD: $r^{\text{target}}_n(s) \leftarrow r(s, a) - (1 - \alpha)\beta \, r(s', a') + \beta W_n(s')$
\STATE Zaktualizuj $W$: $W_{n+1}(s) \leftarrow W_n(s) + \eta_n \left[ \frac{\phi_s(a|s)}{\nu(a|s)} r^{\text{target}}_n(s) - W_n(s) \right]$
\STATE Zaktualizuj $J$: $J_{n+1}(s) \leftarrow J_n(s) + \theta_n \left[ \frac{\mu(a|s)}{\nu(a|s)} r^{\text{target}}_n(s) - J_n(s) \right]$
\ENDFOR
\ENDFOR
\STATE \textbf{Zwraca:} $J_n$ zbieżne do $J^{\alpha, \beta}_{\mu, \phi_s}$; $W_n$ zbieżne do $J^{\alpha, \beta}_{\phi_s}$
\end{algorithmic}
\end{algorithm}
\end{frame}

\begin{frame}{Algorytm 1: Kluczowe mechanizmy}
\textbf{Cel TD (różnica temporalna):}
\[
r^{\text{target}}_n(s) = r(s,a) - (1-\alpha)\beta \, r(s',a') + \beta W_n(s')
\]

\begin{itemize}
    \item Człon $-(1-\alpha)\beta r(s',a')$ -- korekta quasi-hiperboliczna
    \item $\beta W_n(s')$ -- przybliżona wartość quasi-hiperboliczna dla $\phi_s$
\end{itemize}

\vspace{0.5em}
\textbf{Dwie skale czasowe:}
\begin{itemize}
    \item \textbf{Szybka skala} ($\eta_n$): aktualizacja $W_n$ (wartość dla $\phi_s$)
    \item \textbf{Wolna skala} ($\theta_n$): aktualizacja $J_n$ (wartość dla $(\mu, \phi_s)$)
    \item Warunek: $\frac{\theta_n}{\eta_n} \to 0$ -- separacja skal
\end{itemize}

\vspace{0.5em}
\textbf{Korekta z wagami wiarygodności:}
\[
\rho^W = \frac{\phi_s(a|s)}{\nu(a|s)}, \quad \rho^J = \frac{\mu(a|s)}{\nu(a|s)}
\]
\end{frame}

\begin{frame}{Twierdzenie o~zbieżności Algorytmu 1}
\begin{block}{Twierdzenie 2 \cite{eshwar2025}}
Jeżeli kroki uczenia $\eta_n$, $\theta_n$ spełniają warunki Robbins--Monro:
\[
\sum_n \eta_n = \infty, \quad \sum_n \eta_n^2 < \infty \quad \text{oraz} \quad \sum_n \theta_n = \infty, \quad \sum_n \theta_n^2 < \infty,
\]
oraz dodatkowo $\frac{\theta_n}{\eta_n} \to 0$ (dwie skale czasowe), to ciągi $(W_n, J_n)$ są zbieżne prawie na pewno do $(J^{\alpha, \beta}_{\phi_s}, J^{\alpha, \beta}_{\mu, \phi_s})$.
\end{block}

\vspace{1em}
\textbf{Dowód:} \cite{eshwar2025,borkar2008}.
\end{frame}

\begin{frame}{Optymalna polityka dla dyskontowania QH}
\textbf{Definicja:} Optymalna polityka w dyskontowaniu QH to para optymalna $(\mu^\star, \phi_s^\star)$ maksymalizująca wartość funkcji:
\[
(\mu^\star, \phi_s^\star) = \arg\max_{\mu \in \Phi_m} \max_{\phi_s \in \Phi_s} J^{\alpha, \beta}_{\mu, \phi_s}(s), \quad \forall s \in S.
\]

\vspace{0.5em}
\textbf{Struktura optymalnej polityki:}
\begin{itemize}
    \item $\phi_s^\star = \phi_\beta^\star$ -- optymalna polityka dla dyskontowania wykładniczego
    \item $\mu^\star(s) = \arg\max_{\mu}\sum\limits_{a}\mu(a\mid s) \left[ r(s, a) + \alpha \beta \sum_{s'} q(s'\mid s, a) J^{\beta}_{\phi_\beta^\star}(s') \right]$
\end{itemize}

\vspace{0.5em}
\begin{block}{Twierdzenie 3 \cite{eshwar2025}}
Istnieje deterministyczna optymalna para $(\mu^\star, \phi_s^\star)$ taka, że
\[
J^{\alpha, \beta}_{\mu^\star, \phi_s^\star}(s) = \max_{\mu} \max_{\phi_s} J^{\alpha, \beta}_{\mu, \phi_s}(s), \quad \forall s \in S.
\]
\end{block}

\vspace{0.5em}
\textbf{Dowód:} 1285\_MS.pdf.
\end{frame}

\begin{frame}{Algorytm 2: QH Q-Learning}
\begin{algorithm}[H]
\caption{Znajdowanie optymalnej polityki dla dyskontowania quasi-hiperbolicznego \cite{eshwar2025}}
\label{alg:qh-qlearning-defense}
\begin{algorithmic}[1]
\small
\STATE \textbf{Wymagania:} Kroki uczenia $(\eta_n)$, $(\theta_n)$; współczynniki dyskontowania $\alpha$, $\beta$
\STATE \textbf{Inicjalizacja:} $W_0, Q_0 \in \mathbb{R}^{|S| \times |A|}$
\FOR{$n = 0, 1, 2, \ldots$}
\FOR{$(s, a) \in S \times A$}
\STATE Próbkuj $s' \sim q(\cdot|s, a)$ i zaobserwuj $r(s, a)$
\STATE $W_{n+1}(s, a) \leftarrow W_n(s, a) + \eta_n \left[ r(s, a) + \beta \max_{a'} W_n(s', a') - W_n(s, a) \right]$
\STATE $Q_{n+1}(s, a) \leftarrow Q_n(s, a) + \theta_n \left[ (1 - \alpha) r(s, a) + \alpha W_n(s, a) - Q_n(s, a) \right]$
\ENDFOR
\ENDFOR
\STATE \textbf{Zwraca:} $Q_n$ (QH) oraz $W_n$ (wykładnicze);
\[ \mu^\star(s) = \arg\max\limits_a Q_n(s, a), \]
\[ \phi_s^\star(s) = \arg\max\limits_a W_n(s, a). \]
\end{algorithmic}
\end{algorithm}

\vspace{0.3em}
\textbf{Kluczowe własności:}
\begin{itemize}
    \item Algorytm bezmodelowy -- nie wymaga znajomości $q$
    \item $W$ uczy się funkcji Q dla dyskontowania wykładniczego
    \item $Q$ uczy się funkcji Q dla dyskontowania quasi-hiperbolicznego
\end{itemize}
\end{frame}

\begin{frame}{Relacja między funkcjami Q}
\textbf{Dla pary polityk $(\mu, \phi_s)$:}
\[
Q^{\alpha, \beta}_{\mu, \phi_s}(s,a) = (1-\alpha) r(s,a) + \alpha Q^{\beta}_{\phi_s}(s,a)
\]

\vspace{0.5em}
\textbf{Dla optymalnych polityk:}
\begin{align*}
Q^{\alpha,\beta}_{\mu^\star,\phi_s^\star}(s,a) &= (1-\alpha)\,r(s,a) + \alpha\,Q^{\beta}_{\phi_s^\star}(s,a) \\
\phi_s^\star &= \arg\max_{\phi_s} Q^{\beta}_{\phi_s}(s,a) = \phi_\beta^\star \\
\mu^\star(s) &= \arg\max_a Q^{\alpha, \beta}_{\mu^\star, \phi_s^\star}(s,a)
\end{align*}

\vspace{0.5em}
\textbf{Aktualizacje w~Algorytmie 2:}
\begin{itemize}
    \item $W$ -- klasyczny Q-learning (wykładniczy)
    \item $Q$ -- liniowa kombinacja nagrody bieżącej i~$W$ (quasi-hiperboliczny)
    \item Separacja skal: $W$ zbiega szybciej, $Q$ traktuje $W$ jako quasi-stałe
\end{itemize}
\end{frame}

\begin{frame}{Twierdzenie o~zbieżności Algorytmu 2}
\begin{block}{Twierdzenie 4 \cite{eshwar2025}}
Załóżmy, że kroki uczenia $(\eta_n)$ i $(\theta_n)$ spełniają warunki Robbins--Monro:
\[
\sum_n \eta_n = \infty, \quad \sum_n \eta_n^2 < \infty, \quad \sum_n \theta_n = \infty, \quad \sum_n \theta_n^2 < \infty,
\]
oraz dodatkowo $\theta_n / \eta_n \to 0$ (dwie skale czasowe).

Wtedy ciągi $(Q_n, W_n)$ generowane przez Algorytm QH Q-Learning zbiegają do $(Q^{\alpha, \beta}_{\mu^\star, \phi_s^\star}, Q^{\beta}_{\phi_s^\star})$.

Para $\mu^\star(s) = \arg\max_a Q_n(s,a)$ oraz $\phi_s^\star(s) = \arg\max_a W_n(s,a)$ jest optymalna.
\end{block}

\vspace{1em}
\textbf{Dowód:} \cite{eshwar2025,borkar2008}.
\end{frame}

% ==========================================
% Eksperymenty
% ==========================================
\section{Eksperymenty i~wyniki}

\begin{frame}{Środowiska testowe}
\textbf{Parametry wspólne:}
\begin{itemize}
    \item Polityka próbkowania: $\nu$ -- rozkład jednostajny (eksploracja)
    \item Kroki uczenia: $\eta_n = \frac{\eta_0}{(c + n)^{\kappa}}$, $\theta_n = \frac{\theta_0}{(c + n)^{\lambda}}$ gdzie $\frac{1}{2} < \kappa < \lambda \le 1$
\end{itemize}

\vspace{0.5em}
\textbf{Trzy środowiska testowe:}
\begin{enumerate}
    \item \textbf{MDP z~dwoma stanami}
    \begin{itemize}
        \item Prosty model do weryfikacji analitycznej
        \item 2 stany, 2 akcje
        \item Możliwość dokładnego obliczenia wartości teoretycznych
    \end{itemize}
    
    \vspace{0.3em}
    \item \textbf{Problem zarządzania zapasami}
    \begin{itemize}
        \item Dyskretny magazyn o~pojemności $M=3$
        \item Losowy popyt, koszty magazynowania i~zamówienia
        \item Test algorytmów w~realistycznym scenariuszu biznesowym
    \end{itemize}
    
    \vspace{0.3em}
    \item \textbf{Problem przetrwania rozbitka}
    \begin{itemize}
        \item Złożony model zarządzania zasobami
        \item 11~616 stanów, 28 akcji, dynamika stochastyczna
        \item Test skalowalności w~dużych przestrzeniach
    \end{itemize}
\end{enumerate}
\end{frame}

\begin{frame}{MDP z~dwoma stanami -- Model}
\begin{columns}[T]
\begin{column}{0.48\textwidth}
\textbf{Parametry:}
\begin{itemize}
    \item Stany: $S = \{1, 2\}$
    \item Akcje: $A = \{a_1, a_2\}$
    \item $\alpha = 0.5$, $\beta = 0.9$
\end{itemize}

\vspace{0.3em}
\textbf{Nagrody:}
\begin{itemize}
    \item $r(1,a_1) = 0$, $r(1,a_2) = 2$
    \item $r(2,a_1) = 15$
    \item $r(2,a_2,1) = 20$, $r(2,a_2,2) = 5$
\end{itemize}

\vspace{0.3em}
\textbf{Testowane polityki:}
\begin{itemize}
    \item 4 deterministyczne
    \item 1 mieszana (50/50)
\end{itemize}
\end{column}

\begin{column}{0.48\textwidth}
\centering
\includegraphics[width=0.95\textwidth]{example1.png}
\small Diagram przejść w~modelu dwustanowym
\end{column}
\end{columns}
\end{frame}

\begin{frame}{MDP z~dwoma stanami -- Wyniki Algorytmu 1}
\textbf{Porównanie wartości teoretycznych i aproksymowanych (po $1\times10^6$ iteracji):}

\begin{table}[h!]
    \centering
    % ZMIANA: Pierwszy parametr to \textwidth (maksymalna szerokość slajdu), 
    % drugi to '!' (wysokość dopasuje się automatycznie, zachowując proporcje).
    \resizebox{\textwidth}{!}{%
    \begin{tabular}{llcrrrr}
        \toprule
        $\mu$, $\phi_s$ & Stan & $W$ (teoria) & $W$ (aproksymata) & $J^{\alpha,\beta}_{\rho}$ (teoria) & $J$ (aproksymata) \\
        \midrule
        $\,\mu$ = $\phi_1$: zawsze $a_1$ & 1 & 71.0526 & 71.0526 & 35.5263 & 35.5175 \\
        $\,\mu$ = $\phi_1$: zawsze $a_1$ & 2 & 78.9474 & 78.9474 & 46.9737 & 46.9643 \\
        $\,\mu$ = $\phi_2$: zawsze $a_2$ & 1 & 67.2500 & 67.2668 & 34.6250 & 34.6340 \\
        $\,\mu$ = $\phi_2$: zawsze $a_2$ & 2 & 77.7500 & 77.7588 & 45.1250 & 45.1254 \\
        $\,\mu$ = $\phi_3$: $a_1$ w $s=1$, $a_2$ w $s=2$ & 1 & 77.5862 & 77.6618 & 38.7931 & 38.7936 \\
        $\,\mu$ = $\phi_3$: $a_1$ w $s=1$, $a_2$ w $s=2$ & 2 & 86.2069 & 86.3167 & 49.3534 & 49.3612 \\
        $\,\mu$ = $\phi_4$: $a_2$ w $s=1$, $a_1$ w $s=2$ & 1 & 60.3448 & 60.4001 & 31.1724 & 31.1910 \\
        $\,\mu$ = $\phi_4$: $a_2$ w $s=1$, $a_1$ w $s=2$ & 2 & 69.3103 & 69.3575 & 42.1552 & 42.1706 \\
        $\,\mu$ = $\phi_5$: losowa 50/50 & 1 & 69.3534 & 69.2913 & 35.1767 & 35.1625 \\
        $\,\mu$ = $\phi_5$: losowa 50/50 & 2 & 78.1466 & 78.0512 & 45.9483 & 45.9297 \\
        \bottomrule
    \end{tabular}%
    }
\end{table}
\end{frame}

\begin{frame}{MDP z~dwoma stanami -- Zbieżność $W$ i $J$}
\begin{figure}[htbp]
    \centering
    \includegraphics[width=0.95\textwidth]{Ex1PLots.png}
    \caption{Zbieżność aproksymacji w eksperymencie 1: błąd normy $\ell_2$ względem wartości teoretycznych dla wybranych polityk.}
\end{figure}
\end{frame}

\begin{frame}{MDP z~dwoma stanami -- Wyniki Algorytmu 2}
\textbf{Uczenie optymalnej polityki (25 000 000 iteracji):}

\begin{table}[h!]
\centering
\caption{Oszacowane wartości $Q(s,a)$ i $W(s,a)$ po zakończeniu uczenia (stany $S=\{1,2\}$).}
\label{tab:two_state_learned_qw}
\begin{minipage}[t]{0.48\textwidth}
\centering
\begin{tabular}{lrr}
	\toprule
	\textbf{Stan $s$} & $Q(s,a_1)$ & $Q(s,a_2)$ \\
\midrule
1 & 38.7921 & \textbf{38.8527} \\
2 & \textbf{49.9129} & 49.3519 \\
\bottomrule
\end{tabular}
\end{minipage}
\hfill
\begin{minipage}[t]{0.48\textwidth}
\centering
\begin{tabular}{lrr}
	\toprule
	\textbf{Stan $s$} & $W(s,a_1)$ & $W(s,a_2)$ \\
\midrule
1 & \textbf{77.5813} & 75.6954 \\
2 & 84.8224 & \textbf{86.1988} \\
\bottomrule
\end{tabular}
\end{minipage}
\end{table}

\vspace{0.3em}
\begin{table}[h!]
\centering
\caption{Wartości analityczne $Q(s,a)$ i $W(s,a)$ dla polityk optymalnych w MDP dwustanowym.}
\label{tab:two_state_analytic_qw}
\begin{minipage}[t]{0.48\textwidth}
\centering
\begin{tabular}{lrr}
    \toprule
    \textbf{Stan $s$} & $Q(s,a_1)$ & $Q(s,a_2)$ \\
\midrule
1 & 38.7931 & \textbf{38.8534} \\
2 & \textbf{49.9138} & 49.3534 \\
\bottomrule
\end{tabular}
\end{minipage}
\hfill
\begin{minipage}[t]{0.48\textwidth}
\centering
\begin{tabular}{lrr}
    \toprule
    \textbf{Stan $s$} & $W(s,a_1)$ & $W(s,a_2)$ \\
\midrule
1 & \textbf{77.5862} & 75.7069 \\
2 & 84.8276 & \textbf{86.2069} \\
\bottomrule
\end{tabular}
\end{minipage}
\end{table}

\vspace{0.3em}
\textbf{Optymalne polityki:}
\begin{itemize}
    \item \textbf{$\mu^\star(1)=a_2$, $\mu^\star(2)=a_1$}
    \item \textbf{$\phi_s^\star(1)=a_1$, $\phi_s^\star(2)=a_2$}
\end{itemize}
\end{frame}

\begin{frame}{Problem zarządzania zapasami}
\textbf{Model:}
\begin{itemize}
    \item Pojemność magazynu: $M = 3$ (stany: $s\in\{0,1,2,3\}$)
    \item Popyt losowy: $D \in \{0,1,2,3\}$ z~rozkładem $(0.1, 0.2, 0.3, 0.4)$
    \item Koszty i cena: $c=6$, $h=3$, $p=11$
    \item Parametry: $\alpha = 0.35$, $\beta = 0.95$
\end{itemize}

\vspace{0.5em}
\textbf{Testowane pary polityk:}
\begin{itemize}
    \item $\rho_1=(\hat{\mu},\hat{\phi}_s)$
    \item $\rho_2=(\nu,\hat{\phi}_s)$
    \item $\rho_3=(\hat{\mu},\nu)$
\end{itemize}
\end{frame}

\begin{frame}{Problem zarządzania zapasami -- Wyniki Algorytmu 2}
\begin{table}[h!]
\centering
\caption{Oszacowane wartości $Q(s,a)$ i $W(s,a)$ po zakończeniu uczenia ($M=3$, $\alpha=0.35$, $\beta=0.95$).}
\label{tab:inv_m3_learned_qw}
\begin{minipage}[t]{0.49\textwidth}
\centering
\resizebox{\textwidth}{!}{%
\begin{tabular}{c|rrrr}
    \toprule
$Q(s,a)$ & $a=0$ & $a=1$ & $a=2$ & $a=3$ \\
\midrule
$s=0$ & 44.6349 & 48.4375 & \textbf{49.8316} & 47.6371 \\
$s=1$ & 54.4310 & \textbf{55.8304} & 53.6284 & 47.6284 \\
$s=2$ & \textbf{61.8311} & 59.6326 & 53.6321 & 47.6324 \\
$s=3$ & \textbf{65.6328} & 59.6270 & 53.6232 & 47.6306 \\
\bottomrule
\end{tabular}
}
\end{minipage}
\hfill
\begin{minipage}[t]{0.49\textwidth}
\centering
\resizebox{\textwidth}{!}{%
\begin{tabular}{c|rrrr}
    \toprule
$W(s,a)$ & $a=0$ & $a=1$ & $a=2$ & $a=3$ \\
\midrule
$s=0$ & 127.4992 & 131.6772 & \textbf{134.2060} & 134.1781 \\
$s=1$ & 137.6581 & 140.1643 & \textbf{140.2175} & 134.2134 \\
$s=2$ & \textbf{146.1939} & 146.1299 & 140.2028 & 134.1548 \\
$s=3$ & \textbf{152.2200} & 146.2367 & 140.1319 & 134.2182 \\
\bottomrule
\end{tabular}
}
\end{minipage}
\end{table}

\vspace{0.3em}
\begin{table}[h!]
\centering
\caption{Wartości analityczne $Q^{\mu^\star,\phi_s^\star}(s,a)$ i $W^{\phi_s^\star}(s,a)$ dla modelu $M=3$.}
\label{tab:inv_m3_analytic_wq}
\begin{minipage}[t]{0.49\textwidth}
\centering
\resizebox{\textwidth}{!}{%
\begin{tabular}{c|rrrr}
    \toprule
$Q(s,a)$ & $a=0$ & $a=1$ & $a=2$ & $a=3$ \\
\midrule
$s=0$ & 44.555 & 48.354 & \textbf{49.753} & 47.550 \\
$s=1$ & 54.354 & \textbf{55.753} & 53.550 & 47.550 \\
$s=2$ & \textbf{61.753} & 59.550 & 53.550 & 47.550 \\
$s=3$ & \textbf{65.550} & 59.550 & 53.550 & 47.550 \\
\bottomrule
\end{tabular}
}
\end{minipage}
\hfill
\begin{minipage}[t]{0.49\textwidth}
\centering
\resizebox{\textwidth}{!}{%
\begin{tabular}{c|rrrr}
    \toprule
$W(s,a)$ & $a=0$ & $a=1$ & $a=2$ & $a=3$ \\
\midrule
$s=0$ & 127.30 & 131.47 & 133.98 & \textbf{134.00} \\
$s=1$ & 137.47 & 139.98 & \textbf{140.00} & 134.00 \\
$s=2$ & 145.98 & \textbf{146.00} & 140.00 & 134.00 \\
$s=3$ & \textbf{152.00} & 146.00 & 140.00 & 134.00 \\
\bottomrule
\end{tabular}
}
\end{minipage}
\end{table}
\end{frame}

\begin{frame}{Problem zarządzania zapasami -- Wyniki Algorytmu 1}

\begin{table}[h!]
    \centering
    \caption{Porównanie wartości $W$ i $J$ -- teoria vs aproksymacja (5 000 000 iteracji).}
    \label{tab:inventory_comparison}
    
    % Skalujemy tabelę do 90% szerokości slajdu, wysokość dopasuje się automatycznie
    \resizebox{0.8\textwidth}{!}{%
    \begin{tabular}{llrrrr}
        \toprule
        \textbf{Polityka} & \textbf{Stan $s$} & \textbf{W (teoria)} & \textbf{W (aproks.)} & \textbf{J (teoria)} & \textbf{J (aproks.)} \\
        \midrule
        $\rho_1=(\mu^\star,\phi^\star)$ & 0 & 47.5500 & 48.8707 & 49.7530 & 49.8370 \\
        $\rho_1=(\mu^\star,\phi^\star)$ & 1 & 53.5500 & 53.4860 & 55.7530 & 55.9488 \\
        $\rho_1=(\mu^\star,\phi^\star)$ & 2 & 59.5500 & 59.6520 & 61.7530 & 61.9084 \\
        $\rho_1=(\mu^\star,\phi^\star)$ & 3 & 65.5500 & 66.0263 & 65.5500 & 65.6837 \\
        \midrule
        $\rho_2=(\nu,\phi^\star)$ & 0 & 47.5500 & 48.4988 & 47.5531 & 47.9197 \\
        $\rho_2=(\nu,\phi^\star)$ & 1 & 53.5500 & 53.6143 & 54.5525 & 54.9220 \\
        $\rho_2=(\nu,\phi^\star)$ & 2 & 59.5500 & 60.0191 & 60.6515 & 61.0314 \\
        $\rho_2=(\nu,\phi^\star)$ & 3 & 65.5500 & 66.2218 & 65.5500 & 65.9100 \\
        \midrule
        $\rho_3=(\mu^\star,\nu)$ & 0 & 35.6724 & 35.6749 & 37.8963 & 37.8967 \\
        $\rho_3=(\mu^\star,\nu)$ & 1 & 42.7300 & 42.8086 & 43.8963 & 43.9461 \\
        $\rho_3=(\mu^\star,\nu)$ & 2 & 48.9185 & 48.8855 & 49.8963 & 49.8965 \\
        $\rho_3=(\mu^\star,\nu)$ & 3 & 53.9407 & 54.0234 & 53.9407 & 53.9650 \\
        \bottomrule
    \end{tabular}%
    }
\end{table}

\end{frame}

\begin{frame}{Problem zarządzania zapasami -- Zbieżność}
\begin{figure}[htbp]
    \centering
    \includegraphics[width=0.8\textwidth]{ex2Plots.png}
    \caption{Zbieżność błędu normy $\ell_2$ aproksymacji $W$ i~$J$ dla trzech par polityk (skala logarytmiczna)}
\end{figure}

\end{frame}

\begin{frame}{Problem przetrwania rozbitka -- Model}
\textbf{Założenia:}
\begin{itemize}
    \item Rozbitek na bezludnej wyspie zarządza zasobami (jedzenie, woda, zdrowie)
    \item Stan: $(f, w, h, d)$ -- jedzenie (0--10), woda (0--10), zdrowie (0--5), dni (0--15)
    \item Przestrzeń stanów: $11 \times 11 \times 6 \times 16 = 11\,616$ stanów
    \item Limit czasu: 15 dni (potem wyspa zalewa się)
\end{itemize}

\vspace{0.3em}
\textbf{Akcje (28 kombinacji):}
\begin{itemize}
    \item Tryb: czekaj, połów, ekspedycja ratunkowa
    \item Konsumpcja jedzenia: 0, 1 lub 2 jednostki
    \item Konsumpcja wody: 0, 1 lub 2 jednostki
\end{itemize}

\vspace{0.3em}
\textbf{Dynamika:}
\begin{itemize}
    \item Pogoda: deszcz z~prawd. 50\% (+2 lub +3 wody)
    \item Połów: wymaga 1 jedz.+1 wody, daje 3 jedzenia (szansa rośnie ze zdrowiem)
    \item Ekspedycja: wymaga 7 jedz.+7 wody oraz $h \geq 3$, szansa zależy od zdrowia i~zapasów
    \item Ratunek losowy: 7\% dziennie od dnia 2
\end{itemize}
\end{frame}

\begin{frame}{Problem przetrwania rozbitka -- Wyniki}
\textbf{Symulacja Monte Carlo (500 000 prób):}

\begin{table}[h!]
\centering
\begin{tabular}{lr}
\toprule
\textbf{Polityka} & \textbf{Wskaźnik przeżycia [\%]} \\
\midrule
Impulsywna ($\mu^\star$) & 41.85 \\
Planowa ($\phi_s^\star$) & 43.26 \\
Mieszana ($\mu^\star$, $\phi_s^\star$) & 43.19 \\
Losowa (jednostajna) & 12.81 \\
Człowiek (50 prób) & 42.00 \\
\bottomrule
\end{tabular}
\end{table}

\end{frame}

\begin{frame}{Problem przetrwania rozbitka -- Analiza}
\begin{columns}[T]
\begin{column}{0.48\textwidth}
\centering
\includegraphics[width=0.95\textwidth]{distributioncastaway.png}
\small Rozkład zgonów i~ratunków
\end{column}
\begin{column}{0.48\textwidth}
\centering
\includegraphics[width=0.95\textwidth]{expeditions.png}
\small Rozkład wypraw ratunkowych
\end{column}
\end{columns}
\end{frame}

% ==========================================
% Podsumowanie
% ==========================================
\section{Podsumowanie}

\begin{frame}{Dalsze kierunki badań i~wyzwania}
\textbf{Dalsze prace:}
\begin{itemize}
    \item Psychologia uzależnień 
    \item Gry wieloagentowe z~dyskontowaniem QH
    \item Deep Q-Learning dla jeszcze większych przestrzeni stanów
    \item Zastosowania w~planowaniu finansowym i~systemach rekomendacji
\end{itemize}

\vspace{0.5em}
\textbf{Wyzwania badawcze:}
\begin{itemize}
    \item \textbf{Estymacja parametrów:} Jak wyznaczać $\alpha$ i~$\beta$ z~danych rzeczywistych?
    \item \textbf{Adaptacja online:} Czy parametry QH mogą się zmieniać w~czasie?
    \item \textbf{Weryfikacja:} Trudność w~testowaniu na rzeczywistych użytkownikach
\end{itemize}
\end{frame}

\begin{frame}[plain]
\begin{center}
{\Huge Dziękuję za uwagę}

\vspace{2em}

\vspace{3em}
\normalsize
Michał Wrona\\
Politechnika Wrocławska\\
Wydział Matematyki\\[0.5em]
Promotor: prof. dr hab. inż. Anna Jaśkiewicz
\end{center}
\end{frame}

% ==========================================
% Bibliografia z pliku zewnętrznego
% ==========================================
\begin{frame}[allowframebreaks]{Literatura}
    % Styl bibliografii (plain = numerowany, abbrv = skrócony)
    \bibliographystyle{plain}
    % Wskazanie pliku .bib (bez rozszerzenia)
    \bibliography{references}
\end{frame}

\end{document}